%!TEX root = ../../main.tex
%% ===============================================================
%% 11.1 결론
%% 파일: chapters/11_conclusion/01_conclusion.tex
%% 상위: chapters/11_conclusion/chapter.tex
%% 정의 라벨: -
%% 참조 라벨: -
%% 포함 객체: -
%% 인용: -
%% 상태: 초안 (docs/OUTLINE.md 참고)
%% ===============================================================

\section{결론}

본 논문은 분 단위 공개 텔레메트리라는 제약 아래에서 리그 오브 레전드 교전 결과 예측 문제를 정식화하고, 이를 위한 완전한 파이프라인을 제안·평가하였다. 킬 조건부 교전 국소화 알고리즘은 사람의 주석 없이 \numMatches{} 경기로부터 약 100 만 건의 교전을 결정론적으로 찾아내며, 교환가치 라벨은 킬·오브젝트·현상금이 섞인 교전의 순이득을 이진 결과로 요약한다. 무누출 계약을 명시한 다중 모달 표현 위에서 25 종 이상의 모델을 동일 조건으로 비교한 결과, 공학적 테이블 요약과 그래디언트 부스팅의 조합이 시험 AUC $\aucLGBM{}$로 가장 우수하였고, 시퀀스·그래프·사건 기반 신경망 표현은 $.569$--$\aucBiGRU{}$에 머물렀다. 정합 입력 비교는 이 격차의 절반이 학습기 계열에서, 나머지 절반이 표현에서 비롯됨을 보였다. 예측 가능성은 후반 국면과 일방적 골드 격차에서 $.8$에 이르렀으나 대등한 교전에서는 우연 수준에 근접하였는데, 이는 관측의 정보 입도가 부과하는 상한으로 해석된다. 연구 질문에 대한 답은 다음과 같이 요약된다.

\begin{description}[leftmargin=3.2em,style=nextline]
  \item[RQ1] 킬 조건부 규칙만으로 재현 가능한 교전 말뭉치를 구축할 수 있다. 사람 주석과의 일치도는 두 트랙의 검증 프로토콜로 계량되며 결과는 진행 중이다.
  \item[RQ2] 공학적 테이블 요약이 가장 많은 신호를 드러낸다. 시퀀스·그래프·사건 표현의 추가는 신호를 더 드러내지 못했다.
  \item[RQ3] 입력을 고정했을 때 그래디언트 부스팅이 MLP보다 $\gapLearner{}$ 높다.
  \item[RQ4] 예측 가능성은 국면·골드 격차에 따라 단조 증가하며, 대등한 교전의 낮은 예측 가능성은 정보 입도 한계로 설명된다.
  \item[RQ5] 일곱 처치의 효과는 절제 프로토콜로 계량된다. \todo{최종 실행 결과 반영.}
\end{description}
